\chapter{Présentation de l'organisme d'accueil et cadrage du stage}
\section{Présentation générale de Smart Automation Technologies}
Smart Automation Technologies (SAT) est une startup technologique située à Tanger. Elle intervient dans les domaines de la transformation digitale, de l'architecture des systèmes d'information, de la data science et de l'intelligence artificielle. Son positionnement est celui d'un partenaire technologique capable d'accompagner les organisations depuis la compréhension du besoin jusqu'à la conception de solutions adaptées.

Ce contexte est cohérent avec le sujet du stage, car la transformation d'un WMS vers un Smart WMS exige d'abord une compréhension métier solide. Il ne suffit pas d'ajouter des technologies à un processus existant ; il faut identifier les problèmes opérationnels, les données disponibles, les dépendances avec les systèmes externes et les conditions réelles d'activation des capacités avancées.

\section{Contexte et périmètre du stage}
Le stage couvre la période d'août à octobre 2026 et porte sur la conception d'une architecture intelligente pour transformer un WMS traditionnel en Smart WMS. Le périmètre est volontairement centré sur la modélisation métier, fonctionnelle et applicative logique. Les choix définitifs de technologie, de base de données, d'infrastructure cloud ou d'équipement physique restent hors périmètre à ce stade.

\section{Objectifs}
Les objectifs opérationnels du stage sont les suivants :
\begin{itemize}[leftmargin=1.2cm]
  \item réaliser un benchmark fonctionnel des solutions WMS leaders ;
  \item identifier les fonctions communes et les limites d'un WMS traditionnel ;
  \item formaliser les use cases Core, Smart et Optional ;
  \item proposer une architecture fonctionnelle cible modulaire ;
  \item préparer une modélisation UML cohérente ;
  \item traduire la cible fonctionnelle vers une architecture applicative logique.
\end{itemize}

\section{Livrables}
Les livrables attendus sont un benchmark WMS, une cartographie des use cases, une architecture fonctionnelle cible, des modèles UML et une proposition d'architecture applicative. Le projet logiciel local SmartWMS AI Operations Agent est traité comme une maquette d'aide à la décision : il démontre certains principes d'outillage, de règles déterministes et d'agent conversationnel, sans prétendre constituer un WMS complet.
